\expandafter\gdef\csname SNside1\endcsname{\node[anchor=north west,inner sep=0pt,text width=48mm,font=\SideFont,align=left] at ([xshift=138mm,yshift=-534.39pt]current page.north west) {\textcolor{SNink}{\textbf{课程连接 · 套接字由操作系统管理}}\par\vspace{1.5mm}应用通过套接字接口使用传输服务，端口参与内核分用，进程号标识操作系统中的进程。一个服务器进程可接受多条连接，故不能把端口、进程和连接按一一对应关系死记。};}
\expandafter\gdef\csname SNside2\endcsname{\node[anchor=north west,inner sep=0pt,text width=48mm,font=\SideFont,align=left] at ([xshift=138mm,yshift=-397.86pt]current page.north west) {\textcolor{SNink}{\textbf{最小反例 · 同端口仍有两条连接}}\par\vspace{1.5mm}客户A:50000与客户B:50000同时连接服务器S:80，服务端目的端口完全相同，但源IP不同，TCP四元组不同。只看服务端端口就合并连接，会把两个客户的序号和状态混成一套。};}
\expandafter\gdef\csname SNside3\endcsname{\node[anchor=north west,inner sep=0pt,text width=48mm,font=\SideFont,align=left] at ([xshift=138mm,yshift=-95.00pt]current page.north west) {\textcolor{SNink}{\textbf{不变量 · UDP长度包含8字节首部}}\par\vspace{1.5mm}应用数据12 B时UDP长度为20 B；普通IPv4还要把IP首部计入整个数据报上限。UDP自身保留报文边界，承载它的IP分片则必须先在目的IP层重组，不能把每个IP片都当一个UDP报文。};}
\expandafter\gdef\csname SNside4\endcsname{\node[anchor=north west,inner sep=0pt,text width=48mm,font=\SideFont,align=left] at ([xshift=138mm,yshift=-95.00pt]current page.north west) {\textcolor{SNink}{\textbf{错项排除 · 伪首部不是真的多发一层}}\par\vspace{1.5mm}校验借用源IP、目的IP、协议和长度，是为检查交付对象；它只参加计算。奇数字节补0也只补计算过程。因此计算线上长度和吞吐量时，不能再给UDP额外增加伪首部和求和补位。};}
\expandafter\gdef\csname SNside5\endcsname{\node[anchor=north west,inner sep=0pt,text width=48mm,font=\SideFont,align=left] at ([xshift=138mm,yshift=-95.00pt]current page.north west) {\textcolor{SNink}{\textbf{不变量 · 回卷在取反之前完成}}\par\vspace{1.5mm}B9B6加65C5得到带进位的11F7B，先回卷成1F7C，再取反成E083。直接丢掉最高进位会得到E084；若未取反则仍是中间和，辨认选项时两处步骤都要检查。};}
\expandafter\gdef\csname SNside6\endcsname{\node[anchor=north west,inner sep=0pt,text width=48mm,font=\SideFont,align=left] at ([xshift=138mm,yshift=-95.00pt]current page.north west) {\textcolor{SNink}{\textbf{最小反例 · 双向seq各有自己的坐标}}\par\vspace{1.5mm}甲已发2001～4000，乙回复seq=4001、ack=3001：乙的seq属于反方向，只确认甲到3000的是ack。把乙的seq拿来确认甲的数据，会错误释放仍在途的3001～4000。};}
\expandafter\gdef\csname SNside7\endcsname{\node[anchor=north west,inner sep=0pt,text width=48mm,font=\SideFont,align=left] at ([xshift=138mm,yshift=-95.00pt]current page.north west) {\textcolor{SNink}{\textbf{课程连接 · 状态是事件触发的记录}}\par\vspace{1.5mm}连接状态可像操作系统进程状态一样按事件画转换：发FIN、收到其ACK、收到对方FIN分别触发不同动作。仅看到报文名字不足以决定两端都已关闭，还要确定报文在何时被哪端收到。};}
\expandafter\gdef\csname SNside8\endcsname{\node[anchor=north west,inner sep=0pt,text width=48mm,font=\SideFont,align=left] at ([xshift=138mm,yshift=-95.00pt]current page.north west) {\textcolor{SNink}{\textbf{最小反例 · 首次加两次还差一次}}\par\vspace{1.5mm}先收到把确认号推进到3001的正常ACK，再收到两份ack=3001，重复计数只有2。第三份重复ACK到达才满足传统快重传阈值；同号ACK共有4份中的第一份不能混算成重复。\par\smallskip\begin{tikzpicture}[x=1mm,y=1mm,>=Stealth,node distance=5mm]
\node[draw=SNink!50,fill=SNblue,rounded corners=1mm,text width=34mm,align=center,inner sep=1.2mm,font=\fontsize{7.5}{9}\selectfont] (m0) at (0,0) {首ACK};
\node[draw=SNink!50,fill=SNblue,rounded corners=1mm,text width=34mm,align=center,inner sep=1.2mm,font=\fontsize{7.5}{9}\selectfont] (m1) at (0,-11) {重复1};
\node[draw=SNink!50,fill=SNblue,rounded corners=1mm,text width=34mm,align=center,inner sep=1.2mm,font=\fontsize{7.5}{9}\selectfont] (m2) at (0,-22) {重复2};
\node[draw=SNink!50,fill=SNblue,rounded corners=1mm,text width=34mm,align=center,inner sep=1.2mm,font=\fontsize{7.5}{9}\selectfont] (m3) at (0,-33) {尚未触发};
\draw[->,SNink] (m0) -- (m1);
\draw[->,SNink] (m1) -- (m2);
\draw[->,SNink] (m2) -- (m3);
\end{tikzpicture}
\par{\fontsize{7.5}{9}\selectfont 三个同号ACK仅含两份重复ACK}};}
\expandafter\gdef\csname SNside9\endcsname{\node[anchor=north west,inner sep=0pt,text width=48mm,font=\SideFont,align=left] at ([xshift=138mm,yshift=-95.00pt]current page.north west) {\textcolor{SNink}{\textbf{最小反例 · 窗口额度不是新发额度}}\par\vspace{1.5mm}总窗口3000 B，其中1000 B已在途，尚可新发只有2000 B。如果再把3000 B全部发出去，未确认量将变成4000 B，越过了当前窗口上限；先扣已占部分才能安排新段。};}
\expandafter\gdef\csname SNside10\endcsname{\node[anchor=north west,inner sep=0pt,text width=48mm,font=\SideFont,align=left] at ([xshift=138mm,yshift=-95.00pt]current page.north west) {\textcolor{SNink}{\textbf{不变量 · 按ACK事件依次更新四账}}\par\vspace{1.5mm}先更新累计确认，再核对已发未确认区间，随后按拥塞算法改变cwnd并读当前rwnd，最后取小扣在途。若只把两个窗口数字取小就回答还能发多少，会在所有有在途数据的题目中多算。};}
\expandafter\gdef\csname SNside11\endcsname{\node[anchor=north west,inner sep=0pt,text width=48mm,font=\SideFont,align=left] at ([xshift=138mm,yshift=-95.00pt]current page.north west) {\textcolor{SNink}{\textbf{错项排除 · 零窗口不等于超时丢包}}\par\vspace{1.5mm}接收进程始终不取数据，rwnd会逐步降到0，即使此前每段都已正确到达。此时暂停普通新数据并运行持续计时器；不能无事件依据把cwnd按超时规则重置，也不能假设探测会腾空缓存。};}
\expandafter\gdef\csname SNside12\endcsname{\node[anchor=north west,inner sep=0pt,text width=48mm,font=\SideFont,align=left] at ([xshift=138mm,yshift=-95.00pt]current page.north west) {\textcolor{SNink}{\textbf{不变量 · 增长轨迹要固定记录时刻}}\par\vspace{1.5mm}某轮发送前cwnd为8，与这一轮ACK处理后cwnd为16是两个时刻。教材的逐RTT数列必须先选定记录点，尤其在跨门限与超时处；初值和门限也必须来自题设，不能默认为现代实现都取1 MSS。};}
\expandafter\gdef\csname SNside13\endcsname{\node[anchor=north west,inner sep=0pt,text width=48mm,font=\SideFont,align=left] at ([xshift=138mm,yshift=-268.78pt]current page.north west) {\textcolor{SNink}{\textbf{错项排除 · 没有拥塞仍可能线性长}}\par\vspace{1.5mm}窗口从8增到32若已处于拥塞避免，每RTT增加1而非翻倍。``未发生拥塞''只排除了丢包反应，没有证明始终在慢开始，也没有排除接收窗口限制，因此必须先判断增长阶段。};}
\expandafter\gdef\csname SNside14\endcsname{\node[anchor=north west,inner sep=0pt,text width=48mm,font=\SideFont,align=left] at ([xshift=138mm,yshift=-95.00pt]current page.north west) {\textcolor{SNink}{\textbf{不变量 · RTT样本先审身份}}\par\vspace{1.5mm}同一字节在0 ms原发、100 ms重发，150 ms收到ACK，样本可能是150 ms也可能是50 ms。没有额外消歧信息时不能随便选一个更新SRTT，Karn原则正是要排除这种歧义。};}
\expandafter\gdef\csname SNside15\endcsname{\node[anchor=north west,inner sep=0pt,text width=48mm,font=\SideFont,align=left] at ([xshift=138mm,yshift=-168.61pt]current page.north west) {\textcolor{SNink}{\textbf{错项排除 · 半关闭仍允许反向数据}}\par\vspace{1.5mm}被动端收到FIN后进入CLOSE-WAIT，表示对方不会再发送，本端应用仍可继续向对方发送。把TCP全双工看成一根双向同时关闭的管道，会错误断言收到FIN后所有方向的数据都不能发。};}
\expandafter\gdef\csname SNside16\endcsname{\node[anchor=north west,inner sep=0pt,text width=48mm,font=\SideFont,align=left] at ([xshift=138mm,yshift=-95.00pt]current page.north west) {\textcolor{SNink}{\textbf{不变量 · 2MSL计在正确的一端}}\par\vspace{1.5mm}正常主动关闭者在收到对方FIN并发最后ACK后进入TIME-WAIT；被动者收到最后ACK即可关闭。求服务器最短释放时间时，应先判服务器是主动还是被动，再决定是否包含这个等待。};}
\AddToShipoutPictureFG{\begin{tikzpicture}[remember picture,overlay]\ifcsname SNside\thepage\endcsname\csname SNside\thepage\endcsname\fi\end{tikzpicture}}